\documentclass{article}
\usepackage[utf8]{inputenc}
\usepackage{amsmath}
\usepackage{graphicx}
\usepackage{hyperref}
\usepackage{microtype}
\usepackage[section]{placeins}
\usepackage{booktabs}
\usepackage{float}
\usepackage[toc,title]{appendix}
\usepackage{listings}
\usepackage{color}
\definecolor{codegreen}{rgb}{0,0.6,0}
\definecolor{codegray}{rgb}{0.5,0.5,0.5}
\definecolor{codepurple}{rgb}{0.58,0,0.82}
\definecolor{backcolour}{rgb}{0.95,0.95,0.92}
\lstdefinestyle{mystyle}{
  backgroundcolor=\color{backcolour},
  commentstyle=\color{codegreen},
  keywordstyle=\color{magenta},
  numberstyle=\tiny\color{codegray},
  stringstyle=\color{codepurple},
  basicstyle=\ttfamily\footnotesize,
  breakatwhitespace=false,
  breaklines=true,
  captionpos=b,
  keepspaces=true,
  numbers=left,
  numbersep=5pt,
  showspaces=false,
  showstringspaces=false,
  showtabs=false,
  tabsize=2
}
\lstset{style=mystyle}
\title{543 Readmissions Project Data Preparation Report}
\author{Jack Phelan}
\date{May 2026}
\begin{document}
\maketitle
\tableofcontents
\newpage

% =============================================================================
\section{Overview}
% Briefly describe the purpose of this report and the starting dataset.
% Input: healthcare_readmissions_dataset_train_post_eda.csv (post-EDA, age outliers already removed)
% Output: healthcare_readmissions_dataset_train_preprocessed.csv
% =============================================================================

\par This report documents the steps taken following exploratory data analysis of
\texttt{healthcare\_readmissions\_dataset\_train.csv}. From EDA, the dataset had 7958 rows
and 19 features (including the target).

% =============================================================================
\section{Data Cleaning}
% =============================================================================

\subsection{Missing Data}
% Describe which variables had missing values and decisions made.
% number_of_prior_visits and medications_prescribed had missings.
% Overlap between the two was small, so row-wise deletion was not used.
% Imputation was deferred until after feature engineering (since medications_prescribed was transformed).
% Final decision: number_of_prior_visits imputed to 0 (see Section 4).
\par Only features \texttt{number\_of\_prior\_visits} and \texttt{medications\_prescribed}
contained missing values (314 (4\%) and 657 (8\%) respectively). Among all rows with missing values,
only 21 lacked both features, so plainly removing all rows with missing values was ruled out.
Imputing strategies were deferred until after feature engineering as transformations were anticipated.

\subsection{Outlier Handling}
% BMI outliers removed using z-score threshold of 3.5.
% This reduced the dataset from 8,038 to 7,951 rows.
% Age outliers were handled in the prior EDA phase.

\par During EDA, 80 rows were identified as outliers and removed for having impossible age values. BMI outliers were noted but not addressed.

\par Using a $z$-score threshold of $|z| > 3.5$, 7 rows were identified for having BMI values outside the plausible range. These rows were removed,
reducing the dataset to 7,951 rows, from 8,038 original and 7,958 post-EDA.

\subsection{Errors and Duplicates}
% Both handled in the EDA phase; nothing additional required here.

\par No additional error correction or duplicate removal was necessary at this stage; both had been addressed during EDA.

% =============================================================================
\section{Feature Engineering}
% =============================================================================

\subsection{Single Variable Transforms}
\label{subsec:single-var}

\subsubsection{\texttt{medications\_prescribed} $\to$ \texttt{is\_prescribed}}
% Binary encoding: 1 if medications_prescribed > 0, else 0.
% Original column dropped. Missing values implicitly mapped to 0.
\par \texttt{medications\_prescribed} was binarised into \texttt{is\_prescribed} (1 if the original value was greater than zero, 0 otherwise),
As there was a long tail of values and for simpler imputation. Missing values were treated as 0.
Other encoding apporaches might be looked at in the future (e.g.\ ordinal bins),
but this was decided to be the baseline.

\subsubsection{\texttt{length\_of\_stay} $\to$ \texttt{length\_of\_stay\_score}}
% Converted to an ordinal score (1–7) following the LACE Index scoring convention.
% LACE: Length of stay, Acuity of admission, Comorbidity, Emergency department use.
% Bins: ≤1 day → 1, 2 days → 2, 3 days → 3, 4–6 days → 4, 7–14 days → 5, >14 days → 7.
% Original column dropped.

\par Continuous length of stay (days) was mapped to an ordinal risk score (1--7) following
the LACE Index convention~\cite{van2010derivation}, a widely used readmission risk scoring system.
The mapping is shown in Table~\ref{tab:los_score}.

\begin{table}[H]
  \centering
  \caption{Length-of-Stay Score Encoding (LACE Index)}
  \label{tab:los_score}
  \begin{tabular}{lc}
    \toprule
    Length of Stay (days) & Score \\
    \midrule
    $\leq 1$    & 1 \\
    2           & 2 \\
    3           & 3 \\
    4--6        & 4 \\
    7--14       & 5 \\
    $> 14$      & 7 \\
    \bottomrule
  \end{tabular}
\end{table}

\noindent The original \texttt{length\_of\_stay} column was dropped after encoding.

\subsubsection{\texttt{age} $\to$ \texttt{age\_group}}
% Binned into clinical risk tiers: 0–18, 19–25, 26–40, 41–65, 66–80, 81+.
% Rationale: age has a non-linear relationship with readmission risk; bins capture clinical tiers.
% Original column dropped.

\par Raw age was binned into rough age groups to capture the non-linear relationship between age and readmission risk, as observed during EDA.

\subsubsection{\texttt{number\_of\_prior\_visits} (type conversion)}
% Cast from object/string to float to allow downstream imputation and modelling.

\par \texttt{number\_of\_prior\_visits} was cast to \texttt{float} to allow downstream imputation and use in numeric models.

% =============================================================================
\section{Feature Selection Decisions}
% =============================================================================

\par The following variables were removed from the dataset:

\begin{itemize}
  \item \textbf{\texttt{weight\_kg}}: dropped because \texttt{adjusted\_weight\_kg} is highly correlated and is presumed to be the clinically adjusted version of the same measurement.
  \item \textbf{\texttt{age}}: replaced by \texttt{age\_group} (see Section~\ref{subsec:single-var}).
  \item \textbf{\texttt{length\_of\_stay}}: replaced by \texttt{length\_of\_stay\_score} (see Section~\ref{subsec:single-var}).
  \item \textbf{\texttt{medications\_prescribed}}: replaced by \texttt{is\_prescribed} (see Section~\ref{subsec:single-var}).
\end{itemize}

\noindent Oversampling and standardization techniques were not applied at this stage,
as more thorough experimentation seems needed to find the best approach for this dataset.

\par Imputing techniques were also kept as simple as possible for this baseline version.
\texttt{number\_of\_prior\_visits} was imputed to 0, but like medication imputing, will be explored further in future iterations.

% =============================================================================
\section{Final Dataset Summary}
\label{sec:dataset-summary}
% Describe the final shape, remaining variables, and any last-minute decisions.
% =============================================================================

\par \texttt{number\_of\_prior\_visits} missing values were imputed to 0 as a conservative default.

\par The final preprocessed dataset contains \textbf{7,951 rows} and the columns listed in Table~\ref{tab:dataset_info}.

\begin{table}
  \caption{Dataset Info}
  \label{tab:dataset_info}
  \begin{tabular}{lrl}
    \toprule
    Column & Non-Null Count & Dtype \\
    \midrule
    patient\_id & 7951 & int64 \\
    gender & 7951 & str \\
    ethnicity & 7951 & str \\
    hospital\_id & 7951 & str \\
    height\_m & 7951 & float64 \\
    smoker & 7951 & bool \\
    bmi & 7951 & float64 \\
    adjusted\_weight\_kg & 7951 & float64 \\
    has\_diabetes & 7951 & int64 \\
    has\_hypertension & 7951 & int64 \\
    exercise\_frequency & 7951 & str \\
    diet\_type & 7951 & str \\
    number\_of\_prior\_visits & 7951 & float64 \\
    type\_of\_treatment & 7951 & str \\
    readmission\_within\_30\_days & 7951 & int64 \\
    is\_prescribed & 7951 & int64 \\
    length\_of\_stay\_score & 7951 & int64 \\
    age\_group & 7951 & category \\
    \bottomrule
  \end{tabular}
\end{table}

\par The file was saved to \texttt{data/interim/healthcare\_readmissions\_dataset\_train\_preprocessed.csv}.

% =============================================================================
\begin{thebibliography}{9}
  \bibitem{van2010derivation}
  van Walraven, C., Dhalla, I.A., Bell, C., et al.\ (2010).
  Derivation and validation of an index to predict early death or unplanned readmission after discharge from hospital to the community.
  \textit{CMAJ}, 182(6), 551--557.
\end{thebibliography}

\end{document}
